% GENERATED FILE. Edit canonical lessons, then run npm run generate:books.
% canonical-source-hash: fnv1a64:e44f48763655a6bd
% canonical-lessons: JA-C162-kowareru, JA-C162-wareru, JA-C162-toreru, JA-C162-shuuri, JA-C162-furui

\chapter{Broken, Cracked and Repaired}
\label{ch:ja-broken-cracked-and-repaired}
\hlchaptermodality{162}

\begin{chapteropening}
\textbf{By the end of this chapter:} \emph{I can say when things break and get repaired.}

Complete the last lesson of chapter 162: I can say when things break and get repaired.
\end{chapteropening}

\section[\texorpdfstring{\ja{こわれる} (kowareru)}{こわれる (kowareru)}]{\ja{こわれる} (kowareru) — to be broken; to break down}

\label{lesson:JA-C162-kowareru}



\begin{quote}
Before the new one: say the Japanese for to get dry, then the Japanese for to dry.
\end{quote}

\subsection*{You'll want to know: \ja{こわれる}}
\textbf{\ja{こわれる}} — \emph{kowareru} — \textquotedblleft{}to be broken; to break down\textquotedblright{}.

\begin{quote}
\textbf{\ja{やすい} \ja{かさは} \ja{よく} \ja{こわれます。}} — \emph{yasui kasa wa yoku kowaremasu} — \textquotedblleft{}Cheap umbrellas often break.\textquotedblright{}
\end{quote}

Polite \textbf{\ja{こわれます}}, \emph{kowaremasu}; an ichidan verb. Its transitive partner is \textbf{\ja{こわす}}, \emph{kowasu}.

\begin{cousinweb}
\textbf{\ja{こわれる}} (\emph{kowareru}) is a native Japanese word (wago), inherited rather than borrowed from Chinese, and built on the same root as \emph{kowasu}, \textquotedblleft{}to break\textquotedblright{}; its kanji was attached later as a kun'yomi (JMdict, KANJIDIC2).
\end{cousinweb}

\begin{grammarlens}[title={What you've built}]
One more word for everyday life at A2.
\end{grammarlens}

\subsection*{Your turn}
\begin{itemize}
\raggedright
  \item \emph{Say it:} \emph{kowareru}
  \item \emph{Say it:} \emph{kowareru}, once more
  \item \emph{Say it:} \emph{yasui kasa wa yoku kowaremasu} — Cheap umbrellas often break.
  \item \emph{Recall:} say the Japanese for to get dry, then the Japanese for to dry, then say \emph{kowareru} again
\end{itemize}

\begin{tcolorbox}[breakable,colback=teal!4,colframe=teal!35!black,title={Before you move on}]
What does \ja{こわれる} mean? (\textquotedblleft{}To be broken; to break down\textquotedblright{}.) Say it once more.
\end{tcolorbox}
\section[\texorpdfstring{\ja{われる} (wareru)}{われる (wareru)}]{\ja{われる} (wareru) — to shatter; to split}

\label{lesson:JA-C162-wareru}



\begin{quote}
Before the new one: say the Japanese for to dry, then the Japanese for to be broken.
\end{quote}

\subsection*{You'll want to know: \ja{われる}}
\textbf{\ja{われる}} — \emph{wareru} — \textquotedblleft{}to shatter; to split\textquotedblright{}.

\begin{quote}
\textbf{\ja{この} \ja{いけの} \ja{こおりは} \ja{よく} \ja{われます。}} — \emph{kono ike no kōri wa yoku waremasu} — \textquotedblleft{}The ice on this pond often cracks.\textquotedblright{}
\end{quote}

Polite \textbf{\ja{われます}}, \emph{waremasu}; an ichidan verb.

\begin{cousinweb}
\textbf{\ja{われる}} (\emph{wareru}) is a native Japanese word (wago), inherited rather than borrowed from Chinese; its kanji was attached later as a kun'yomi, and its older origin is not settled (JMdict, KANJIDIC2).
\end{cousinweb}

\begin{grammarlens}[title={What you've built}]
One more word for everyday life at A2.
\end{grammarlens}

\subsection*{Your turn}
\begin{itemize}
\raggedright
  \item \emph{Say it:} \emph{wareru}
  \item \emph{Say it:} \emph{wareru}, once more
  \item \emph{Say it:} \emph{kono ike no kōri wa yoku waremasu} — The ice on this pond often cracks.
  \item \emph{Recall:} say the Japanese for to dry, then the Japanese for to be broken, then say \emph{wareru} again
\end{itemize}

\begin{tcolorbox}[breakable,colback=teal!4,colframe=teal!35!black,title={Before you move on}]
What does \ja{われる} mean? (\textquotedblleft{}To shatter; to split\textquotedblright{}.) Say it once more.
\end{tcolorbox}
\section[\texorpdfstring{\ja{とれる} (toreru)}{とれる (toreru)}]{\ja{とれる} (toreru) — to come off; to come loose}

\label{lesson:JA-C162-toreru}



\begin{quote}
Before the new one: say the Japanese for to be broken, then the Japanese for to shatter.
\end{quote}

\subsection*{You'll want to know: \ja{とれる}}
\textbf{\ja{とれる}} — \emph{toreru} — \textquotedblleft{}to come off; to come loose\textquotedblright{}.

\begin{quote}
\textbf{\ja{この} \ja{はこの} \ja{ふたは} \ja{よく} \ja{とれます。}} — \emph{kono hako no futa wa yoku toremasu} — \textquotedblleft{}The lid of this box often comes off.\textquotedblright{}
\end{quote}

Polite \textbf{\ja{とれます}}, \emph{toremasu}; an ichidan verb.

\begin{cousinweb}
\textbf{\ja{とれる}} (\emph{toreru}) is a native Japanese word (wago), inherited rather than borrowed from Chinese, and built on the same root as \emph{toru}, \textquotedblleft{}to take\textquotedblright{}; its kanji was attached later as a kun'yomi (JMdict, KANJIDIC2).
\end{cousinweb}

\begin{grammarlens}[title={What you've built}]
One more word for everyday life at A2.
\end{grammarlens}

\subsection*{Your turn}
\begin{itemize}
\raggedright
  \item \emph{Say it:} \emph{toreru}
  \item \emph{Say it:} \emph{toreru}, once more
  \item \emph{Say it:} \emph{kono hako no futa wa yoku toremasu} — The lid of this box often comes off.
  \item \emph{Recall:} say the Japanese for to be broken, then the Japanese for to shatter, then say \emph{toreru} again
\end{itemize}

\begin{tcolorbox}[breakable,colback=teal!4,colframe=teal!35!black,title={Before you move on}]
What does \ja{とれる} mean? (\textquotedblleft{}To come off; to come loose\textquotedblright{}.) Say it once more.
\end{tcolorbox}
\section[\texorpdfstring{\ja{しゅうり} (shūri)}{しゅうり (shūri)}]{\ja{しゅうり} (shūri) — repair; mending}

\label{lesson:JA-C162-shuuri}



\begin{quote}
Before the new one: say the Japanese for to shatter, then the Japanese for to come off.
\end{quote}

\subsection*{You'll want to know: \ja{しゅうり}}
\textbf{\ja{しゅうり}} — \emph{shūri} — \textquotedblleft{}repair; mending\textquotedblright{}.

\begin{quote}
\textbf{\ja{ちちは} \ja{いすの} \ja{しゅうりが} \ja{とくいです。}} — \emph{chichi wa isu no shūri ga tokui desu} — \textquotedblleft{}My father is good at mending chairs.\textquotedblright{}
\end{quote}

\begin{cousinweb}
\textbf{\ja{しゅうり}} (\emph{shūri}) is Sino-Japanese (kango): two Chinese characters read with their on'yomi, readings borrowed from Chinese. They are \emph{shū}, \textquotedblleft{}discipline, master\textquotedblright{} (Middle Chinese \emph{sjuw}) and \emph{ri}, \textquotedblleft{}reason, arrangement\textquotedblright{} (\emph{liX}) (KANJIDIC2; Middle Chinese in Baxter's transcription).
\end{cousinweb}

\begin{grammarlens}[title={What you've built}]
One more word for everyday life at A2.
\end{grammarlens}

\subsection*{Your turn}
\begin{itemize}
\raggedright
  \item \emph{Say it:} \emph{shūri}
  \item \emph{Say it:} \emph{shūri}, once more
  \item \emph{Say it:} \emph{chichi wa isu no shūri ga tokui desu} — My father is good at mending chairs.
  \item \emph{Recall:} say the Japanese for to shatter, then the Japanese for to come off, then say \emph{shūri} again
\end{itemize}

\begin{tcolorbox}[breakable,colback=teal!4,colframe=teal!35!black,title={Before you move on}]
What does \ja{しゅうり} mean? (\textquotedblleft{}Repair; mending\textquotedblright{}.) Say it once more.
\end{tcolorbox}
\section[\texorpdfstring{\ja{ふるい} (furui)}{ふるい (furui)}]{\ja{ふるい} (furui) — old (of things)}

\label{lesson:JA-C162-furui}



\begin{quote}
Before the new one: say the Japanese for to come off, then the Japanese for repair.
\end{quote}

\subsection*{You'll want to know: \ja{ふるい}}
\textbf{\ja{ふるい}} — \emph{furui} — \textquotedblleft{}old (of things)\textquotedblright{}.

\begin{quote}
\textbf{\ja{この} \ja{いすは} \ja{ふるいです。}} — \emph{kono isu wa furui desu} — \textquotedblleft{}This chair is old.\textquotedblright{}
\end{quote}

\begin{cousinweb}
\textbf{\ja{ふるい}} (\emph{furui}) is a native Japanese word (wago), inherited rather than borrowed from Chinese; its kanji was attached later as a kun'yomi, and its older origin is not settled (JMdict, KANJIDIC2).
\end{cousinweb}

\begin{grammarlens}[title={What you've built}]
One more word for everyday life at A2.
\end{grammarlens}

\subsection*{Your turn}
\begin{itemize}
\raggedright
  \item \emph{Say it:} \emph{furui}
  \item \emph{Say it:} \emph{furui}, once more
  \item \emph{Say it:} \emph{kono isu wa furui desu} — This chair is old.
  \item \emph{Recall:} say the Japanese for to come off, then the Japanese for repair, then say \emph{furui} again
\end{itemize}

\begin{tcolorbox}[breakable,colback=teal!4,colframe=teal!35!black,title={Before you move on}]
What does \ja{ふるい} mean? (\textquotedblleft{}Old (of things)\textquotedblright{}.) Say it once more.
\end{tcolorbox}
